\pg{354} II. Dada a existencia de um só ente \emph{composto}, deve admitir-se a existencia de um Ente absolutamente \emph{simples} (¹).

\textsc{Artigo} IV

\emph{Ente immutavel e ente mutavel}.

\textbf{§1. Ente immutavel é ente mutavel.} --- Ente \emph{immutavel} é o que não está sujeito á mudança, isto é, que não pode passar de um estado para outro, da potencia para o acto. --- Ente \emph{mutavel} é o que está sujeito á mudança, isto é, que pode passar de um estado para outro, da potencia para o acto\footnote{O homem, sendo um ente \emph{composto}, conhece primeiramente os entes \emph{compostos}, e por estes os \emph{simples}. Por isso, se não existisse nenhum ente \emph{composto}, o homem não poderia conhecer nenhum ente \emph{simples}, que não é objecto. proporcionado para nós. Mas, se existir um só ente \emph{composto}, podemos e devemos concluir que existe o Ente absolutamente \emph{simples}; porque o \emph{composto}, como dissemos, não pode ser o ente primeiro, mas suppõe e exige uma causa extrinseca, que, em ultima analyse, deve ser simples.}.

\textbf{§2. Especies da immutabilidade.} --- A \emph{immutabilidade} pode ser \emph{absoluta} e \emph{relativa}. --- É \emph{absoluta}, quando o ente exclue toda e qualquer mudança, quer esta derive de um principio intrinseco, quer derive de um principio extrinseco; tal é a \emph{immutabilidade} de \textsc{Deus}. --- É \emph{relativa}, quando o ente só exclue uma ou outra especie de mudança; tal é a \emph{immutabilidade} das creaturas, que, se são immutaveis quanto á \emph{essencia}, não são immutaveis quanto á \emph{existencia}\footnote{A \emph{immutabilidade} é uma consequencia da simplicidade, como a \emph{mutabilidade} o é da composição. A mudança importa uma passagem da potencia para o acto; por isso, onde se encontra a potencia, ahi se encontra a possibilidade da mudança, --- e, onde não se encontra a potencia, ahi se encontra a immutabilidade. --- A \emph{mudança} é, como dissemos, a passagem de um estado para outro. A \emph{mutabilidade} é a disposição do ente para passar de um estado para outro. Diz S. Thomaz: «De ratione \emph{mutationis} est quod aliquid idem se habeat \emph{aliter nunc} et \emph{prius}» (\emph{Sum}. \emph{Th}., p. I, q. 45, a. 2 ad 2).}.

\textbf{§3. Elementos da mudança.} --- Toda a mudança, por ser uma passagem da potencia para o acto, de um opposto para\footnote{Toda a mudança deriva --- ou de um principio intrinseco ao sujeito, que se muda, --- ou de um principio extrinseco, de que o sujeito mutavel depende. Se, portanto, um ente é \emph{immutavel} emquanto exclue toda e qualquer potencialidade, tanto em relação a si mesmo, quanto em ordem ás causas extrinsecas, a sua immutabilidade é \emph{absoluta}. --- Se o sujeito, considerado sob um certo aspecto, exclue a possibilidade da mudança, mas, considerado sob um outro aspecto, encerra essa possibilidade da mudança, a sua immutabilidade é \emph{relativa}.}